\documentclass[11pt,a4paper]{article}
% preamble.tex — shared setup for human-verification documents
% Compile each document with:  tectonic <File>.tex

\usepackage[T1]{fontenc}
\usepackage[utf8]{inputenc}
\usepackage{lmodern}
\usepackage{amssymb,amsmath,amsthm}
\usepackage{mathtools}
\usepackage[dvipsnames]{xcolor}
\usepackage{listings}
\usepackage{hyperref}
\usepackage[margin=2.5cm]{geometry}
\usepackage{parskip}

% ── Colours ──────────────────────────────────────────────────────────
\definecolor{leanBg}{HTML}{F7F7F0}
\definecolor{haskellBg}{HTML}{F0F4F7}
\definecolor{keyword}{HTML}{7B3294}
\definecolor{comment}{HTML}{5A8C32}
\definecolor{noteFg}{HTML}{B35900}

% ── Lean 4 listings language ─────────────────────────────────────────
% NOTE: Lean listing blocks use ASCII approximations of unicode symbols
% (e.g. -> for arrow, o. for compose).  The literate key substitutes
% these with proper math-mode renderings in the PDF.
\lstdefinelanguage{Lean4}{
  morekeywords={def,theorem,structure,where,namespace,end,import,fun,
                let,have,show,by,exact,intro,cases,rfl,subst,obtain,
                funext,simp,rw,apply,variable,Type,Prop,instance,
                section,open,class,extends,deriving,sorry,noncomputable},
  sensitive=true,
  morecomment=[l]{--},
  morecomment=[n]{/-}{-/},
  morestring=[b]",
}

% ── Lean listing style ───────────────────────────────────────────────
\lstdefinestyle{lean}{
  language=Lean4,
  backgroundcolor=\color{leanBg},
  basicstyle=\small\ttfamily,
  keywordstyle=\color{keyword}\bfseries,
  commentstyle=\color{comment}\itshape,
  breaklines=true,
  frame=single,
  framerule=0.4pt,
  xleftmargin=4pt,
  xrightmargin=4pt,
  aboveskip=6pt,
  belowskip=4pt,
  literate=
    {->}{{$\to$}}2
    {=>}{{$\Rightarrow$}}2
    {<|}{{$\langle$}}1
    {|>}{{$\rangle$}}1
    {==>}{{$\Longrightarrow$}}3
    { o. }{{$\;\circ\;$}}3
    {forall}{{$\forall$}}1
    {:=}{{$\,:=\,$}}2
    {!=}{{$\neq$}}2
    {<=}{{$\leq$}}2
    {>=}{{$\geq$}}2,
}

% ── Haskell listing style ───────────────────────────────────────────
\lstdefinestyle{haskell}{
  language=Haskell,
  backgroundcolor=\color{haskellBg},
  basicstyle=\small\ttfamily,
  keywordstyle=\color{keyword}\bfseries,
  commentstyle=\color{comment}\itshape,
  breaklines=true,
  frame=single,
  framerule=0.4pt,
  xleftmargin=4pt,
  xrightmargin=4pt,
  aboveskip=6pt,
  belowskip=4pt,
  literate=
    {->}{{$\to$}}2
    {=>}{{$\Rightarrow$}}2,
}

% ── Theorem environments ────────────────────────────────────────────
\theoremstyle{definition}
\newtheorem{defn}{Definition}[section]
\newtheorem{thm}[defn]{Theorem}
\newtheorem{prop}[defn]{Proposition}

\theoremstyle{remark}
\newtheorem*{fidelitynote}{Fidelity note}

% ── Convenience macros ───────────────────────────────────────────────
\newcommand{\lean}[1]{\texttt{#1}}      % inline Lean identifier
\newcommand{\hask}[1]{\texttt{#1}}      % inline Haskell identifier
\newcommand{\Ext}{\mathrm{Ext}}
\newcommand{\Pos}{\mathrm{Pos}}
\newcommand{\Shape}{\mathrm{Shape}}


\title{Human Verification: \texttt{Containers.ComonadConverse}\\[4pt]
  \large From Container Comonads Back to Directed Containers (M2b)}
\author{Machine-generated from \texttt{lean/ComonadConverse.lean} for review by Neil Ghani}
\date{July 2026}

\begin{document}
\maketitle

\begin{abstract}
This document presents every definition and theorem \emph{statement} from the
Lean~4 module \lean{Containers.ComonadConverse} in three parallel forms: the
original Lean, a readable Haskell rendering, and precise mathematical English.
Proofs are omitted --- they are machine-checked by Lean.  The purpose is to let
a human verifier confirm that the formalisation captures the intended
mathematics.
\end{abstract}

\tableofcontents
\bigskip

%===================================================================
\section{Overview}
%===================================================================

This file proves the \textbf{converse} of the Ahman--Chapman--Uustalu
equivalence (\emph{When is a container a comonad?}, LMCS 10(3):2014).  The
forward direction (in \lean{Containers.Directed}) shows that a directed
container induces a comonad on its extension.  Here we run the construction
backwards: from a comonad structure on $\mathbf{Ext}\,C$ we \textbf{recover} a
directed container and prove that the recovered structure satisfies D1--D5.

\medskip
\noindent\textbf{Strategy.}  By the representation theorem
(\lean{Containers.Basic}), any natural transformation between container
extensions is a container morphism.  So the counit and comultiplication of a
comonad on $\mathbf{Ext}\,C$ are determined by combinatorial data:

\begin{align*}
  \varepsilon\langle s, v \rangle &= v(e(s)), & &\text{a root choice }e,\\
  \delta\langle s, v \rangle &=
    \bigl\langle \mathtt{cShape}(s),\;
      \lambda p.\langle \mathtt{cSub}(s,p),\; \lambda q.\; v(\mathtt{cVal}(s,p,q)) \rangle
    \bigr\rangle,
    & &\text{three components.}
\end{align*}

The \textbf{linchpin}: the right counit law forces $\mathtt{cShape} = \mathrm{id}$
(\lean{cShape\_eq}).  Once the outer shape collapses, $(e, \mathtt{cSub}, \mathtt{cVal})$
are exactly a root, sub-shape, and shift, and the three comonad laws become D1--D5 ---
the same partition as the forward direction, read in reverse.

Everything is \textbf{Type}-level, Lean~4 core, no Mathlib.

%===================================================================
\section{Auxiliary lemmas}
%===================================================================

% --- 1. Ext.norm_aux ----------------------------------------------
\subsection{Lemma: \lean{Ext.norm\_aux}}

\begin{prop}[Shape-collapse normalisation]
Let $C$ be a container, $X$ a type, $s, t : C.\mathrm{Shape}$, and $h_c : t = s$.
Then for any $\mathtt{sub} : C.\mathrm{Pos}(t) \to C.\mathrm{Shape}$ and
$\mathtt{val} : (p : C.\mathrm{Pos}(t)) \to C.\mathrm{Pos}(\mathtt{sub}(p)) \to C.\mathrm{Pos}(s)$
and $v : C.\mathrm{Pos}(s) \to X$,
\[
  \bigl\langle t,\; \lambda p.\langle \mathtt{sub}(p),\;
    \lambda q.\; v(\mathtt{val}(p, q))\rangle \bigr\rangle
  \;=\;
  \bigl\langle s,\; \lambda p.\langle \mathtt{sub}(h_c^{-1}{}_* p),\;
    \lambda q.\; v(\mathtt{val}(h_c^{-1}{}_* p, q))\rangle \bigr\rangle
\]
in $\mathbf{Ext}\,C\,(\mathbf{Ext}\,C\,X)$.  This absorbs every
$\mathtt{cShape}$-transport in \lean{comult\_norm}.
\end{prop}

\paragraph{Lean.}
\begin{lstlisting}[style=lean]
theorem Ext.norm_aux {C : Container} {X : Type} {s t : C.Shape} (hc : t = s)
    (sub : C.Pos t -> C.Shape) (val : (p : C.Pos t) -> C.Pos (sub p) -> C.Pos s)
    (v : C.Pos s -> X) :
    (<|t, fun p => <|sub p, fun q => v (val p q)|>|> : Ext C (Ext C X))
      = <|s, fun p => <|sub (hc.symm cast p), fun q => v (val (hc.symm cast p) q)|>|>
\end{lstlisting}

\paragraph{Haskell.}
\begin{lstlisting}[style=haskell]
-- When cShape s = s (propositionally), we can normalise the outer shape.
-- In Haskell this is trivial because there are no dependent transports;
-- the equality would be definitional.
\end{lstlisting}


% --- 2. Ext.map_mk ------------------------------------------------
\subsection{Lemma: \lean{Ext.map\_mk}}

\begin{prop}[Map on a literal pair]
For $f : X \to Y$, shape $s$, and labelling $g : C.\mathrm{Pos}(s) \to X$,
\[
  \mathbf{Ext}.\mathtt{map}\,f\,\langle s, g \rangle
  \;=\;
  \langle s,\; f \circ g \rangle.
\]
This fires only once the argument is in pair form, enabling normalisation to
rewrite the inner comultiplication before unfolding $\mathbf{Ext}.\mathtt{map}$.
\end{prop}

\paragraph{Lean.}
\begin{lstlisting}[style=lean]
theorem Ext.map_mk {C : Container} {X Y : Type} (f : X -> Y) (s : C.Shape)
    (g : C.Pos s -> X) : Ext.map f (<|s, g|> : Ext C X) = <|s, fun i => f (g i)|>
\end{lstlisting}

\paragraph{Haskell.}
\begin{lstlisting}[style=haskell]
-- fmap f (Ext s g) = Ext s (f . g)
-- This is just the Functor instance for Ext, applied to a concrete pair.
\end{lstlisting}


% --- 3. Ext.mk_inj ------------------------------------------------
\subsection{Lemma: \lean{Ext.mk\_inj}}

\begin{prop}[Injectivity of $\mathbf{Ext}$ pairs]
If $\langle s_1, v_1 \rangle = \langle s_2, v_2 \rangle$ in $\mathbf{Ext}\,C\,X$,
then there exists a shape equality $h_s : s_1 = s_2$ such that
$\forall q : C.\mathrm{Pos}(s_1),\; v_1(q) = v_2((h_s)_* q)$.
\end{prop}

\paragraph{Lean.}
\begin{lstlisting}[style=lean]
theorem Ext.mk_inj {C : Container} {X : Type} {s1 s2 : C.Shape}
    {v1 : C.Pos s1 -> X} {v2 : C.Pos s2 -> X}
    (h : (<|s1, v1|> : Ext C X) = <|s2, v2|>) :
    exists hs : s1 = s2, forall q : C.Pos s1, v1 q = v2 (hs cast q)
\end{lstlisting}

\paragraph{Haskell.}
\begin{lstlisting}[style=haskell]
-- If two Ext values are equal, their shapes are equal and their
-- labellings agree (up to coercion on positions).
-- In Haskell: if Ext s1 v1 == Ext s2 v2 then s1 == s2 and v1 = v2 . coerce.
\end{lstlisting}

\begin{fidelitynote}
This is the ``destructor'' dual to \lean{Ext.ext\_eq}: where \lean{ext\_eq}
\emph{builds} an equality from pointwise data, \lean{mk\_inj} \emph{extracts}
data from an equality.  It is used throughout to read the directed-container
laws (D1--D5) out of the comonad-law equalities.
\end{fidelitynote}


%===================================================================
\section{Container comonad data}
%===================================================================

% --- 4. ContainerComonad ------------------------------------------
\subsection{Structure: \lean{ContainerComonad}}

\begin{defn}[Container comonad]
A \textbf{container comonad} $M$ extends a container with the combinatorial
data guaranteed by the representation theorem:
\begin{itemize}
  \item $e : (s : S) \to P(s)$ --- the counit's \emph{root choice},
  \item $\mathtt{cShape} : S \to S$ --- the outer shape of the comultiplication,
  \item $\mathtt{cSub} : (s : S) \to P(\mathtt{cShape}(s)) \to S$ ---
    the inner shapes (future sub-shapes),
  \item $\mathtt{cVal} : (s : S) \to (p : P(\mathtt{cShape}(s))) \to
    P(\mathtt{cSub}(s, p)) \to P(s)$ --- the values (future shift).
\end{itemize}
The counit and comultiplication are then:
\[
  \varepsilon\langle s, v \rangle = v(e(s)),
\]
\[
  \delta\langle s, v \rangle =
  \bigl\langle \mathtt{cShape}(s),\;
    \lambda p.\langle \mathtt{cSub}(s,p),\;
      \lambda q.\; v(\mathtt{cVal}(s,p,q))\rangle
  \bigr\rangle.
\]
\end{defn}

\paragraph{Lean.}
\begin{lstlisting}[style=lean]
structure ContainerComonad extends Container where
  e : (s : Shape) -> Pos s
  cShape : Shape -> Shape
  cSub : (s : Shape) -> Pos (cShape s) -> Shape
  cVal : (s : Shape) -> (p : Pos (cShape s)) -> Pos (cSub s p) -> Pos s
\end{lstlisting}

\paragraph{Haskell.}
\begin{lstlisting}[style=haskell]
data ContainerComonad s p = ContainerComonad
  { shape   :: s -> s          -- outer shape of comult (forced to id)
  , cSub    :: s -> p -> s     -- inner shapes
  , cVal    :: s -> p -> p -> p -- future shift
  , e       :: s -> p           -- root (counit data)
  }
\end{lstlisting}


% --- 5. ContainerComonad.counit -----------------------------------
\subsection{Definition: \lean{ContainerComonad.counit}}

\begin{defn}[Counit]
\[
  \varepsilon : \mathbf{Ext}\,M\,X \to X, \qquad
  \varepsilon\langle s, v \rangle = v(e(s)).
\]
\end{defn}

\paragraph{Lean.}
\begin{lstlisting}[style=lean]
def counit : Ext M.toContainer X -> X :=
  fun t => t.2 (M.e t.1)
\end{lstlisting}

\paragraph{Haskell.}
\begin{lstlisting}[style=haskell]
counit :: ContainerComonad s p -> Ext s p x -> x
counit m (Ext s v) = v (e m s)
\end{lstlisting}


% --- 6. ContainerComonad.comult -----------------------------------
\subsection{Definition: \lean{ContainerComonad.comult}}

\begin{defn}[Comultiplication]
\[
  \delta : \mathbf{Ext}\,M\,X \to \mathbf{Ext}\,M\,(\mathbf{Ext}\,M\,X),
\]
\[
  \delta\langle s, v \rangle =
  \bigl\langle \mathtt{cShape}(s),\;
    \lambda p.\langle \mathtt{cSub}(s,p),\;
      \lambda q.\; v(\mathtt{cVal}(s,p,q))\rangle
  \bigr\rangle.
\]
\end{defn}

\paragraph{Lean.}
\begin{lstlisting}[style=lean]
def comult : Ext M.toContainer X -> Ext M.toContainer (Ext M.toContainer X) :=
  fun t => <|M.cShape t.1, fun p => <|M.cSub t.1 p, fun q => t.2 (M.cVal t.1 p q)|>|>
\end{lstlisting}

\paragraph{Haskell.}
\begin{lstlisting}[style=haskell]
comult :: ContainerComonad s p -> Ext s p x -> Ext s p (Ext s p x)
comult m (Ext s v) = Ext (shape m s)
  (\p -> Ext (cSub m s p) (\q -> v (cVal m s p q)))
\end{lstlisting}


% --- 7. IsComonad -------------------------------------------------
\subsection{Structure: \lean{IsComonad}}

\begin{defn}[Comonad laws]
A container comonad $M$ \textbf{is a comonad} if the following three laws hold
for all types $X$:
\begin{itemize}
  \item \textbf{Left counit:} $\varepsilon \circ \delta = \mathrm{id}$,
  \item \textbf{Right counit:}
    $\mathbf{Ext}.\mathtt{map}(\varepsilon) \circ \delta = \mathrm{id}$,
  \item \textbf{Coassociativity:}
    $\mathbf{Ext}.\mathtt{map}(\delta) \circ \delta = \delta \circ \delta$.
\end{itemize}
\end{defn}

\paragraph{Lean.}
\begin{lstlisting}[style=lean]
structure IsComonad (M : ContainerComonad) : Prop where
  left_counit : forall {X : Type},
    M.counit o. M.comult = (id : Ext M.toContainer X -> Ext M.toContainer X)
  right_counit : forall {X : Type},
    Ext.map M.counit o. M.comult = (id : Ext M.toContainer X -> Ext M.toContainer X)
  coassoc : forall {X : Type},
    Ext.map M.comult o. M.comult
      = (M.comult o. M.comult :
          Ext M.toContainer X ->
            Ext M.toContainer (Ext M.toContainer (Ext M.toContainer X)))
\end{lstlisting}

\paragraph{Haskell.}
\begin{lstlisting}[style=haskell]
-- Comonad laws:
-- left_counit  : extract . comult = id
-- right_counit : fmap extract . comult = id
-- coassoc      : fmap comult . comult = comult . comult
\end{lstlisting}


%===================================================================
\section{Recovering the directed-container operations}
%===================================================================

% --- 8. cShape_eq -------------------------------------------------
\subsection{Theorem: \lean{ContainerComonad.cShape\_eq}}

\begin{thm}[The linchpin: $\mathtt{cShape} = \mathrm{id}$]
If $M$ is a comonad then $\mathtt{cShape}(s) = s$ for every shape $s$.

\emph{Proof sketch.}  Probe the right counit law
$\mathbf{Ext}.\mathtt{map}(\varepsilon) \circ \delta = \mathrm{id}$
at the generic element $\langle s, \mathrm{id} \rangle : \mathbf{Ext}\,M\,(P(s))$
and read off the shape component.
\end{thm}

\paragraph{Lean.}
\begin{lstlisting}[style=lean]
theorem cShape_eq (h : M.IsComonad) (s : M.Shape) : M.cShape s = s
\end{lstlisting}

\paragraph{Haskell.}
\begin{lstlisting}[style=haskell]
-- The right counit law forces the outer shape map to be the identity.
-- Once cShape = id, (e, cSub, cVal) are exactly (root, sub, shift).
\end{lstlisting}

\begin{fidelitynote}
This is the key step of the converse construction.  Every dependent cast in the
rest of the file is a transport along $\mathtt{cShape\_eq}$.  The Lean proof
uses $\mathtt{congrArg}\,\Sigma.\mathtt{fst}$ to project the shape component
from the equality \lean{hr}.
\end{fidelitynote}


% --- 9. rRoot -----------------------------------------------------
\subsection{Definition: \lean{ContainerComonad.rRoot}}

\begin{defn}[Recovered root]
\[
  o^\mathrm{rec}(s) \;:=\; e(s) \;\in\; P(s).
\]
\end{defn}

\paragraph{Lean.}
\begin{lstlisting}[style=lean]
def rRoot (_h : M.IsComonad) (s : M.Shape) : M.Pos s := M.e s
\end{lstlisting}

\paragraph{Haskell.}
\begin{lstlisting}[style=haskell]
rRoot :: ContainerComonad s p -> IsComonad -> s -> p
rRoot m _ s = e m s
\end{lstlisting}


% --- 10. rSub -----------------------------------------------------
\subsection{Definition: \lean{ContainerComonad.rSub}}

\begin{defn}[Recovered sub-shape]
\[
  s \mathop{\downarrow^\mathrm{rec}} p
  \;:=\;
  \mathtt{cSub}\bigl(s,\; (\mathtt{cShape\_eq}(s))^{-1}_* p\bigr),
\]
transporting $p : P(s)$ into the $\mathtt{cShape}(s)$ fibre along
$\mathtt{cShape\_eq}$.
\end{defn}

\paragraph{Lean.}
\begin{lstlisting}[style=lean]
def rSub (s : M.Shape) (p : M.Pos s) : M.Shape :=
  M.cSub s ((M.cShape_eq h s).symm cast p)
\end{lstlisting}

\paragraph{Haskell.}
\begin{lstlisting}[style=haskell]
rSub :: ContainerComonad s p -> IsComonad -> s -> p -> s
rSub m h s p = cSub m s (coerce p)  -- coerce along cShape_eq s
\end{lstlisting}


% --- 11. rShift ---------------------------------------------------
\subsection{Definition: \lean{ContainerComonad.rShift}}

\begin{defn}[Recovered shift]
\[
  p \oplus^\mathrm{rec} q
  \;:=\;
  \mathtt{cVal}\bigl(s,\; (\mathtt{cShape\_eq}(s))^{-1}_* p,\; q\bigr),
\]
transporting $p$ into the $\mathtt{cShape}(s)$ fibre along $\mathtt{cShape\_eq}$.
\end{defn}

\paragraph{Lean.}
\begin{lstlisting}[style=lean]
def rShift (s : M.Shape) (p : M.Pos s) (q : M.Pos (M.rSub h s p)) : M.Pos s :=
  M.cVal s ((M.cShape_eq h s).symm cast p) q
\end{lstlisting}

\paragraph{Haskell.}
\begin{lstlisting}[style=haskell]
rShift :: ContainerComonad s p -> IsComonad -> s -> p -> p -> p
rShift m h s p q = cVal m s (coerce p) q  -- coerce p along cShape_eq s
\end{lstlisting}


% --- 12. comult_norm ----------------------------------------------
\subsection{Theorem: \lean{ContainerComonad.comult\_norm}}

\begin{thm}[Normalised comultiplication]
For any $s : M.\mathrm{Shape}$ and $v : P(s) \to X$,
\[
  \delta\langle s, v \rangle
  \;=\;
  \bigl\langle s,\;
    \lambda p.\langle s \mathop{\downarrow^\mathrm{rec}} p,\;
      \lambda q.\; v(p \oplus^\mathrm{rec} q)\rangle
  \bigr\rangle.
\]
The outer shape $\mathtt{cShape}(s)$ has been collapsed to $s$ via
$\mathtt{cShape\_eq}$.  All downstream uses of the comonad laws yield D1--D5
with only the intrinsic transports, exactly mirroring the forward direction.
\end{thm}

\paragraph{Lean.}
\begin{lstlisting}[style=lean]
theorem comult_norm (s : M.Shape) (v : M.Pos s -> X) :
    M.comult (<|s, v|> : Ext M.toContainer X)
      = <|s, fun p => <|M.rSub h s p, fun q => v (M.rShift h s p q)|>|>
\end{lstlisting}

\paragraph{Haskell.}
\begin{lstlisting}[style=haskell]
-- After normalisation (cShape = id), comult looks exactly like the
-- directed-container comultiplication:
-- comult (Ext s v) = Ext s (\p -> Ext (rSub s p) (\q -> v (rShift s p q)))
\end{lstlisting}


%===================================================================
\section{Recovering D1--D5}
%===================================================================

% --- 13. rD1 ------------------------------------------------------
\subsection{Theorem: \lean{ContainerComonad.rD1}}

\begin{thm}[D1: sub-shape at the root]
\[
  s \mathop{\downarrow^\mathrm{rec}} o^\mathrm{rec}(s) \;=\; s.
\]
Derived from the left counit law, probing at $\langle s, \mathrm{id}\rangle$.
\end{thm}

\paragraph{Lean.}
\begin{lstlisting}[style=lean]
theorem rD1 (s : M.Shape) : M.rSub h s (M.rRoot h s) = s
\end{lstlisting}

\paragraph{Haskell.}
\begin{lstlisting}[style=haskell]
-- D1:  rSub s (rRoot s) = s
\end{lstlisting}


% --- 14. rD2 ------------------------------------------------------
\subsection{Theorem: \lean{ContainerComonad.rD2}}

\begin{thm}[D2: shift by the root is the identity]
For $q : P(s \mathop{\downarrow^\mathrm{rec}} o^\mathrm{rec}(s))$,
\[
  o^\mathrm{rec}(s) \oplus^\mathrm{rec} q
  \;=\;
  (\mathtt{rD1}(s))_* q.
\]
The value half of the left counit law.
\end{thm}

\paragraph{Lean.}
\begin{lstlisting}[style=lean]
theorem rD2 (s : M.Shape) (q : M.Pos (M.rSub h s (M.rRoot h s))) :
    M.rShift h s (M.rRoot h s) q = (M.rD1 h s) cast q
\end{lstlisting}

\paragraph{Haskell.}
\begin{lstlisting}[style=haskell]
-- D2:  rShift s (rRoot s) q = coerce q  (coerce along rD1 s)
\end{lstlisting}


% --- 15. rD3 ------------------------------------------------------
\subsection{Theorem: \lean{ContainerComonad.rD3}}

\begin{thm}[D3: shift by the sub-shape's root is the identity]
For $p : P(s)$,
\[
  p \oplus^\mathrm{rec} o^\mathrm{rec}(s \mathop{\downarrow^\mathrm{rec}} p) \;=\; p.
\]
The value half of the right counit law.
\end{thm}

\paragraph{Lean.}
\begin{lstlisting}[style=lean]
theorem rD3 (s : M.Shape) (p : M.Pos s) :
    M.rShift h s p (M.rRoot h (M.rSub h s p)) = p
\end{lstlisting}

\paragraph{Haskell.}
\begin{lstlisting}[style=haskell]
-- D3:  rShift s p (rRoot (rSub s p)) = p
\end{lstlisting}


% --- 16. rD4 ------------------------------------------------------
\subsection{Theorem: \lean{ContainerComonad.rD4}}

\begin{thm}[D4: sub-shapes nest along the shift]
For $p : P(s)$ and $q : P(s \mathop{\downarrow^\mathrm{rec}} p)$,
\[
  s \mathop{\downarrow^\mathrm{rec}} (p \oplus^\mathrm{rec} q)
  \;=\;
  (s \mathop{\downarrow^\mathrm{rec}} p) \mathop{\downarrow^\mathrm{rec}} q.
\]
The shape half of coassociativity.
\end{thm}

\paragraph{Lean.}
\begin{lstlisting}[style=lean]
theorem rD4 (s : M.Shape) (p : M.Pos s) (q : M.Pos (M.rSub h s p)) :
    M.rSub h s (M.rShift h s p q) = M.rSub h (M.rSub h s p) q
\end{lstlisting}

\paragraph{Haskell.}
\begin{lstlisting}[style=haskell]
-- D4:  rSub s (rShift s p q) = rSub (rSub s p) q
\end{lstlisting}


% --- 17. rD5 ------------------------------------------------------
\subsection{Theorem: \lean{ContainerComonad.rD5}}

\begin{thm}[D5: associativity of shift]
For $p : P(s)$, $q : P(s \mathop{\downarrow^\mathrm{rec}} p)$,
$q' : P((s \mathop{\downarrow^\mathrm{rec}} p) \mathop{\downarrow^\mathrm{rec}} q)$,
\[
  (p \oplus^\mathrm{rec} q) \oplus^\mathrm{rec} (\mathtt{rD4}(s,p,q)^{-1}_* q')
  \;=\;
  p \oplus^\mathrm{rec} (q \oplus^\mathrm{rec} q').
\]
The value half of coassociativity.
\end{thm}

\paragraph{Lean.}
\begin{lstlisting}[style=lean]
theorem rD5 (s : M.Shape) (p : M.Pos s) (q : M.Pos (M.rSub h s p))
    (q' : M.Pos (M.rSub h (M.rSub h s p) q)) :
    M.rShift h s (M.rShift h s p q) ((M.rD4 h s p q).symm cast q')
      = M.rShift h s p (M.rShift h (M.rSub h s p) q q')
\end{lstlisting}

\paragraph{Haskell.}
\begin{lstlisting}[style=haskell]
-- D5 (shifted associativity):
-- rShift s (rShift s p q) (coerce q') = rShift s p (rShift (rSub s p) q q')
-- where coerce is along (rD4 s p q).symm
\end{lstlisting}


%===================================================================
\section{The full equivalence}
%===================================================================

% --- 18. toDirectedContainer --------------------------------------
\subsection{Definition: \lean{ContainerComonad.toDirectedContainer}}

\begin{defn}[Recovery construction (M2b)]
Given a container comonad $M$ and a proof $h : \mathtt{IsComonad}(M)$, the
\textbf{recovered directed container} is
\[
  (M.\mathrm{toContainer},\;
   \mathtt{rRoot}_h,\;
   \mathtt{rSub}_h,\;
   \mathtt{rShift}_h)
\]
with laws \lean{rD1}--\lean{rD5}.
\end{defn}

\paragraph{Lean.}
\begin{lstlisting}[style=lean]
def toDirectedContainer (h : M.IsComonad) : DirectedContainer where
  toContainer := M.toContainer
  root := M.rRoot h
  sub  := M.rSub h
  shift := M.rShift h
  d1 := M.rD1 h
  d4 := M.rD4 h
  d2 := M.rD2 h
  d3 := M.rD3 h
  d5 := M.rD5 h
\end{lstlisting}

\paragraph{Haskell.}
\begin{lstlisting}[style=haskell]
toDirectedContainer :: ContainerComonad s p -> IsComonad -> DirectedContainer s p
toDirectedContainer m h = DirectedContainer
  { root  = rRoot m h
  , sub   = rSub  m h
  , shift = rShift m h
  -- laws D1--D5 proved from the comonad laws
  }
\end{lstlisting}


% --- 19. toDirectedContainer_counit -------------------------------
\subsection{Theorem: \lean{ContainerComonad.toDirectedContainer\_counit}}

\begin{thm}[Counit round-trip]
The recovered directed container induces the \emph{same} counit as the
comonad we started from, on the nose:
\[
  \varepsilon_{\mathtt{toDirectedContainer}(M,h)}
  \;=\;
  \varepsilon_M.
\]
This holds definitionally (\lean{rfl}).
\end{thm}

\paragraph{Lean.}
\begin{lstlisting}[style=lean]
theorem toDirectedContainer_counit (h : M.IsComonad) {X : Type}
    (t : Ext M.toContainer X) :
    (M.toDirectedContainer h).counit (X := X) t = M.counit t
\end{lstlisting}

\paragraph{Haskell.}
\begin{lstlisting}[style=haskell]
-- The recovered counit equals the original counit:
-- counit (toDirectedContainer m h) = counit m
\end{lstlisting}


% --- 20. toDirectedContainer_comult -------------------------------
\subsection{Theorem: \lean{ContainerComonad.toDirectedContainer\_comult}}

\begin{thm}[Comultiplication round-trip]
The recovered directed container induces the \emph{same} comultiplication as
the comonad we started from:
\[
  \delta_{\mathtt{toDirectedContainer}(M,h)}(t)
  \;=\;
  \delta_M(t).
\]
This is \lean{comult\_norm} in reverse, and together with
\lean{toDirectedContainer\_counit} it says: applying the forward construction
(\lean{Containers.Directed}) to the recovered directed container returns the
original comonad.  This is one half of the equivalence
$\mathbf{DirectedContainer} \simeq \text{container comonad}$.
\end{thm}

\paragraph{Lean.}
\begin{lstlisting}[style=lean]
theorem toDirectedContainer_comult (h : M.IsComonad) {X : Type}
    (t : Ext M.toContainer X) :
    (M.toDirectedContainer h).comult (X := X) t = M.comult t
\end{lstlisting}

\paragraph{Haskell.}
\begin{lstlisting}[style=haskell]
-- The recovered comultiplication equals the original:
-- comult (toDirectedContainer m h) = comult m
\end{lstlisting}


%===================================================================
\section{The other direction}
%===================================================================

% --- 21. toContainerComonad ---------------------------------------
\subsection{Definition: \lean{DirectedContainer.toContainerComonad}}

\begin{defn}[Packaging a directed container as comonad data]
Given a directed container $D$, its \textbf{packaged comonad data} is:
\[
  M_D = (D.\mathrm{toContainer},\;
         e := D.\mathrm{root},\;
         \mathtt{cShape} := \mathrm{id},\;
         \mathtt{cSub} := D.\mathrm{sub},\;
         \mathtt{cVal} := D.\mathrm{shift}).
\]
\end{defn}

\paragraph{Lean.}
\begin{lstlisting}[style=lean]
def toContainerComonad (D : DirectedContainer) : ContainerComonad where
  toContainer := D.toContainer
  e := D.root
  cShape := id
  cSub := D.sub
  cVal := D.shift
\end{lstlisting}

\paragraph{Haskell.}
\begin{lstlisting}[style=haskell]
toContainerComonad :: DirectedContainer s p -> ContainerComonad s p
toContainerComonad d = ContainerComonad
  { shape = id, cSub = sub d, cVal = shift d, e = root d }
\end{lstlisting}


% --- 22. toContainerComonad_isComonad -----------------------------
\subsection{Theorem: \lean{DirectedContainer.toContainerComonad\_isComonad}}

\begin{thm}[Packaged directed container is a comonad]
For any directed container $D$, the packaged data $M_D$ satisfies the comonad
laws (these are exactly the forward laws proved in \lean{Containers.Directed}).
\end{thm}

\paragraph{Lean.}
\begin{lstlisting}[style=lean]
theorem toContainerComonad_isComonad (D : DirectedContainer) :
    D.toContainerComonad.IsComonad where
  left_counit := fun {_} => D.left_counit
  right_counit := fun {_} => D.right_counit
  coassoc := fun {_} => D.coassoc
\end{lstlisting}

\paragraph{Haskell.}
\begin{lstlisting}[style=haskell]
-- Directed container laws D1--D5 imply the comonad laws for the packaged data.
\end{lstlisting}


% --- 23. toContainerComonad_toDirectedContainer -------------------
\subsection{Theorem: \lean{DirectedContainer.toContainerComonad\_toDirectedContainer}}

\begin{thm}[The other half of the equivalence]
For any directed container $D$,
\[
  \mathtt{toDirectedContainer}(M_D, h_D) \;=\; D
\]
definitionally (where $h_D$ is the comonad-law proof above).  This holds
because $\mathtt{cShape} = \mathrm{id}$ in $M_D$, so every
$\mathtt{cShape\_eq}$-transport is \lean{rfl} by proof irrelevance.

Together with Theorems~\ref{defn:counit-roundtrip} and~\ref{defn:comult-roundtrip},
this establishes the full equivalence:
\[
  \mathbf{DirectedContainer} \;\simeq\; \text{container comonad}.
\]
\end{thm}

\paragraph{Lean.}
\begin{lstlisting}[style=lean]
theorem toContainerComonad_toDirectedContainer (D : DirectedContainer) :
    D.toContainerComonad.toDirectedContainer
      D.toContainerComonad_isComonad = D
\end{lstlisting}

\paragraph{Haskell.}
\begin{lstlisting}[style=haskell]
-- Round-trip: toDirectedContainer (toContainerComonad d) = d
-- This is the definitional equality half of the equivalence.
\end{lstlisting}

\begin{fidelitynote}
The Lean proof is \lean{rfl}: the two sides are \emph{definitionally} equal,
not just propositionally.  This works because $\mathtt{cShape} = \mathrm{id}$
definitionally in $M_D$, so the recovered $\mathtt{rSub}$ and $\mathtt{rShift}$
are definitionally equal to $D.\mathrm{sub}$ and $D.\mathrm{shift}$.  Proof
irrelevance then collapses the law fields.  This is a satisfying closure of the
equivalence.
\end{fidelitynote}


%===================================================================
\section{Translation conventions}
%===================================================================

\begin{enumerate}
\item \textbf{Representation theorem.}  The whole file rests on the fact
  (proved in \lean{Containers.Basic}) that natural transformations between
  container extensions correspond bijectively to container morphisms.  This
  lets us describe the comonad structure combinatorially in terms of
  $(e, \mathtt{cShape}, \mathtt{cSub}, \mathtt{cVal})$ rather than abstractly
  as endofunctors.

\item \textbf{Generic element probing.}  Each directed-container law is
  \emph{read off} the corresponding comonad law by instantiating at the
  \emph{generic element} $\langle s, \mathrm{id}\rangle \in \mathbf{Ext}\,M\,(P(s))$
  --- a Yoneda-style argument.

\item \textbf{Dependent transports.}  Because $\mathtt{cShape}(s) = s$ is only
  \emph{propositionally} true (not definitionally) in a general
  \lean{ContainerComonad}, transports along $\mathtt{cShape\_eq}$ appear
  throughout.  The auxiliary \lean{Ext.norm\_aux} absorbs all such casts in a
  single rewrite (\lean{comult\_norm}), after which D1--D5 are obtained with
  only the intrinsic transports from D1 and D4 --- matching the forward
  direction exactly.

\item \textbf{Haskell coercion.}  In the Haskell renderings, all transports
  are glossed as \hask{coerce} with a comment indicating the equality they
  transport along.  This is because GHC's \hask{Data.Coerce} provides a
  safe coercion for representationally equal types, but it cannot capture
  the full dependent-type content.

\item \textbf{Proof bodies omitted.}  Only definitions and theorem
  \emph{statements} are translated.  Lean checks the proofs; the human
  verifier checks that the \emph{statements} match the intended mathematics.
\end{enumerate}

\end{document}
